\subsection{Bird's-Eye View LiDAR Representation and Pillar Encoding}
\label{sec:rel_bev_encoding}

Dense horizontal projection provides an efficient pathway to apply high-throughput 2D convolutional networks to unstructured 3D LiDAR point clouds. PIXOR~\cite{yang2018pixor} established the foundational paradigm of proposal-free Bird's-Eye View (BEV) detection by discretizing point clouds into fixed vertical occupancy slices and calibrated reflectance maps. Building upon this formulation, Nguyen et al.~\cite{nguyen2024indoor} introduced MobilePIXOR, replacing heavyweight convolutional layers with inverted residual blocks and Gaussian heatmap supervision for indoor robotic perception. In parallel, PointPillars~\cite{lang2019pointpillars} eliminated rigid vertical slicing by organizing point clouds into vertical columns (pillars) and learning spatial feature representations via localized PointNet encoders. Modern extensions such as PillarNeXt~\cite{li2023pillarnext} and PillarNeSt~\cite{mao2023pillarnest} systematically modernized pillar-based receptive field designs and backbone scaling for large-scale outdoor perception.

Despite their throughput advantages, existing BEV representations exhibit distinct structural trade-offs. On one hand, PointNet-based pillar encoders~\cite{lang2019pointpillars,li2023pillarnext} achieve expressive geometric learning, yet the non-deterministic point gathering, dynamic memory allocation, and sparse scatter operations incur prohibitive latency and energy penalties on edge accelerators. On the other hand, fixed-grid rasterization methods~\cite{yang2018pixor,nguyen2024indoor} bypass point gathering entirely but suffer from severe spatial quantization error on standard $0.1\,\text{m}$ grids; by treating cells as binary flags or raw counters, they discard critical continuous subcell offsets and high-order vertical distributions, causing boundary staircasing and localization jitter. While recent studies such as PillarHist~\cite{zhou2025pillarhist} investigate quantization-aware height histograms and TriBand-BEV~\cite{khoshkdahan2026triband} explores multi-band vertical slicing to assist pedestrian detection, a shared limitation across these paradigms is the absence of continuous subcell centroid displacements coupled with continuous vertical statistical moments inside a compact, hardware-friendly tensor. In summary, retaining fine continuous geometric boundaries without learned PointNet overhead remains an open research gap. MobilePIXORNeXt bridges this gap through the Strata-BEV14 representation (Section~\ref{sec:strata_bev14}). By deterministically integrating eight stratified vertical density bins, four continuous elevation moments ($\mu_z, \sigma_z, z_{\min}, z_{\max}$), and two continuous subcell centroid displacement offsets ($\Delta x_c, \Delta y_c$) into a 14-channel layout, Strata-BEV14 preserves sub-voxel geometric fidelity while reducing Host-to-Device transfer payloads by $60\%$ without learned encoder latency.

\subsection{Lightweight Convolutional Backbones and Efficient Spatial Attention}
\label{sec:rel_backbone_attention}

Extracting discriminative spatial semantics from sparse BEV feature planes requires balancing expressive receptive fields with strict parameter and latency constraints. In 3D LiDAR perception, volumetric detectors such as SECOND~\cite{yan2018second} popularized sparse convolutions to process non-empty voxels, while two-stage point architectures such as PointRCNN~\cite{shi2019pointrcnn} and hybrid frameworks such as Part-$A^2$~\cite{shi2020part} and PV-RCNN~\cite{shi2020pvrcnn} introduced part-aware aggregation and voxel-set abstractions. Although these methods achieve high detection accuracy on benchmarks such as KITTI~\cite{geiger2012kitti}, their multi-million parameter footprints ($4.0\text{--}63.8\,\text{M}$) and heavy memory access latencies ($130\text{--}195\,\text{ms}$) prohibit real-time deployment on power-constrained edge accelerators.

To overcome the latency bottleneck of volumetric processing, lightweight 2D convolutional designs have gained prominence. The ConvNeXt paradigm~\cite{liu2022convnext} demonstrated that conventional convolutions modernized with $7 \times 7$ depthwise kernels, inverted bottleneck expansions ($e=2.5$), and LayerScale can achieve representational power competitive with Vision Transformers while preserving convolutional hardware efficiency. Concurrently, lightweight attention mechanisms have been developed to emphasize salient geometric features. While Efficient Channel Attention (ECA-Net)~\cite{wang2020eca} captures 1D cross-channel interactions with low complexity, it lacks spatial feature modulation. To address this weakness, Yang et al.~\cite{yang2021simam} proposed SimAM, deriving full 3D spatial-channel attention weights analytically via an energy function inspired by visual neuroscience, requiring zero trainable parameters.

Collectively, a shared limitation of prior mobile LiDAR detectors is the reliance on standard inverted residual bottlenecks with small $3 \times 3$ kernels~\cite{nguyen2024indoor}, which struggle to capture multi-scale context across disparate vehicle and pedestrian dimensions. Furthermore, attempting to compensate by appending external attention modules or distillation schemes at coarse downsampling boundaries~\cite{tran2026rubber} introduces architectural overhead without resolving intra-block receptive field bottlenecks under sparse point distributions. In contrast, MobilePIXORNeXt incorporates the MobilePixorNeXt-S backbone with embedded parameter-free Block-SimAM (Sections~\ref{sec:backbone_architecture}--\ref{sec:next_s_block}) and a Top-Down Scale-Gated FPN (Section~\ref{sec:neck_architecture}), capturing expansive $7 \times 7$ depthwise receptive fields and modulating residual feature updates directly inside every residual block with zero added attention parameters.

\subsection{Geometric Supervision, Loss Formulations, and Quality Calibration}
\label{sec:rel_supervision_losses}

Bounding box regression in oriented object detection has progressively evolved from unconstrained coordinate offsets toward geometrically bounded objectives. Early detectors primarily relied on decoupled Smooth-L1 loss functions. However, regressing spatial coordinates, dimensions, and angles independently ignores the geometric interdependence of bounding boxes and induces severe gradient instability. To provide scale-invariant geometric feedback, Rezatofighi et al.~\cite{rezatofighi2019giou} introduced Generalized Intersection over Union (GIoU) to address gradient vanishing when candidate bounding boxes do not overlap. More recently, Le et al.~\cite{le2026mgiou} formulated Marginalized Generalized Intersection over Union (MGIoU), establishing a unified projection-based framework to optimize convex parametric shapes and arbitrary oriented bounding boxes without rotational singularities.

In parallel with geometric loss formulation, joint quality calibration has emerged as an essential mechanism for accurate detection ranking. Li et al.~\cite{li2020gfl} introduced Generalized Focal Loss (GFL), demonstrating that uncalibrated classification scores frequently disagree with actual localization quality, causing low-quality boxes with high confidence to suppress well-localized candidates during Non-Maximum Suppression (NMS). Additionally, Kendall et al.~\cite{kendall2018uncertainty} introduced multi-task uncertainty weighting, utilizing homoscedastic task variance to balance heterogeneous loss objectives across classification and geometric regression tasks.

Despite these advancements, existing oriented bounding box detectors still struggle with three pervasive limitations. First, independent angle and corner regression objectives remain susceptible to the $\pi$-periodicity boundary jump problem, where slight angular deviations around $\pm \pi/2$ or $\pm \pi$ trigger massive loss spikes and yaw gradient collapse. Second, conventional quality calibration mechanisms rely on auxiliary decoupled prediction heads that inflate model parameters and exacerbate gradient contention across heterogeneous object classes. Third, unconstrained multi-task uncertainty weighting allows task variance parameters to diverge, leading to training instability when balancing disparate loss scales across classification, geometric overlap, and orientation. Overall, developing a unified training formulation that simultaneously enforces orientation-symmetric geometric overlap, calibrated quality scoring (with or without dedicated heads), and bounded multi-task loss balancing remains an under-explored research gap. MobilePIXORNeXt bridges this gap through Oriented Geometric Alignment (OGA), its head-free alternative Q-OGA, and IoU-Aware Quality Alignment (IQA) (Section~\ref{sec:losses_and_oga}). Specifically, OGA unifies multi-axis MGIoU projections with scale-normalized $\pi$-symmetric corner supervision to eliminate boundary discontinuities; Q-OGA couples continuous IoU quality directly into classification supervision to achieve quality calibration with zero head overhead; IQA provides explicit multi-branch quality alignment scored via $S_{\mathrm{det}} = \sqrt{P_{\mathrm{cls}}\hat{q}}$ to suppress miscalibrated false positives; and Temperature-Softmax Bounded Uncertainty Weighting (T-SBUW) guarantees stable convergence by constraining multi-task weights within a normalized, bounded simplex.
